\section{Introduction}\label{sec:intro}

Oral diseases, including ulcerative, inflammatory, and
neoplastic lesions, present a complex visual challenge for automated
analysis. Conditions such as recurrent aphthous ulcer, oral lichen planus, and
squamous-cell carcinoma can share overlapping colour, texture, and boundary
cues at the image level, yet a single misclassification between a benign
inflammatory lesion and an early malignancy carries materially different
clinical consequences. A computer-aided diagnostic system must therefore solve
two related but distinct problems: identifying \emph{whether} a lesion is
malignant (binary screening) and identifying \emph{which} of several
visually similar conditions is present (multi-class subtyping). The second
problem is the one that most directly governs the choice of downstream
management pathway, and it is the focus of this paper.

Deep learning has matured rapidly in this domain. CNNs and, more recently,
Vision Transformers have produced strong binary screening
results~\cite{ref5,ref16,ref17} and several promising multi-class
systems~\cite{ref19,ref22}. Despite this progress, three persistent gaps remain
in the literature.

\emph{Gap~1: Binary formulations dominate.} Most published architectures
collapse oral pathology to a benign-versus-malignant decision. This formulation
is clinically insufficient when downstream treatment depends on differentiating
among multiple non-malignant conditions that visually mimic malignancy, since
ulcerative, inflammatory, and neoplastic lesions warrant different management
pathways.

\emph{Gap~2: Attention modules are stacked without ablation.} Modern
backbones routinely embed attention components such as Squeeze-and-Excitation
(SE), the Convolutional Block Attention Module (CBAM), the Bottleneck Attention
Module (BAM), Triplet Attention, Efficient Multi-scale Attention (EMA), and
Kolmogorov-Arnold Network attention (KAN), yet their individual and combined
contributions in the oral-disease setting are rarely measured. The literature
offers little guidance on which attention modules complement one another and
which conflict, and proposed architectures tend to inherit attention designs
from generic benchmarks rather than testing them on the target distribution.

\emph{Gap~3: Explainability evidence is sparse.} Few multi-class oral
disease studies present visual attention maps that demonstrate the model
attends to lesion tissue rather than to teeth, lip outlines, or lighting
artifacts. Absent such evidence, it is difficult to argue that a high-accuracy
model is also clinically defensible.

The present work addresses these three gaps directly. We make the following
four contributions.

\begin{enumerate}
  \item We design \textbf{Custom EfficientNet~V2}, a five-stage
  parameter-efficient backbone derived from EfficientNetV2-B0 in which the
  standard Block-4 is replaced by a domain-tuned \emph{AttentionHub} and the
  redundant Block-6 and convolutional head are removed, yielding a
  $4.79$~M-parameter, $0.493$~GFLOPs model.
  \item We perform a systematic \textbf{seven-cell ablation} of the
  AttentionHub spanning the \{BAM, Triplet, KAN\} singletons, all pairs, and
  the full triple, plus a no-attention control. From this ablation we derive a
  \textbf{role-complementarity principle}: pairing two spatial-role attention
  modules regresses the hub below its no-attention control, whereas pairing a
  spatial module with a purely channel-wise one is complementary. The principle
  is supported by the v1 ablation grid and an independent negative-result run.
  \item Guided by the principle, we introduce \textbf{AttentionHub-v2}, a
  sequential Triplet\arrowto{}SE cascade in which spatial cross-dimensional
  attention is followed by purely channel-wise recalibration.
  AttentionHub-v2 attains the highest subtype accuracy in the entire study
  ($99.51\,\%$) at a parameter and compute budget ($4.79$~M / $0.493$~GFLOPs)
  essentially identical to the v1 hub and well below every baseline.
  \item We benchmark the proposed model against nine standard backbones under
  a \textbf{single matched training recipe}, with all models trained from
  scratch (no ImageNet pre-training), and provide Grad-CAM++ and LIME
  explainability panels for every model. The full ablation results and
  explainability artifacts are released alongside the paper.
\end{enumerate}

Figure~\ref{fig:arch} previews the proposed architecture: a five-stage backbone
whose fourth stage is the AttentionHub and whose pooled feature feeds two
task-specific heads, one for the binary task and one for the seven-class
subtype task. The architectural changes are deliberately minimal (a single
swapped stage and a trimmed tail) so that the measured gains can be
attributed to the hub design rather than to a wholesale redesign of the
network.

The remainder of this paper is organised as follows.
Section~\ref{sec:related} reviews prior work on oral lesion classification and
segmentation, makes the three motivating gaps explicit, and positions the
present study. Section~\ref{sec:prelim} fixes the notation and the attention
role taxonomy on which the ablation rests. Section~\ref{sec:methods} describes
the dataset, the dual-head classifier, the Custom EfficientNet~V2 backbone, the
two AttentionHub variants, the multi-task loss, and the matched training
protocol. Section~\ref{sec:setup} specifies the experimental setup,
reproducibility provisions, and ethical considerations.
Section~\ref{sec:results} reports the baseline benchmark and per-class
results. Section~\ref{sec:ablation} presents the ablation study and develops
the role-complementarity principle. Section~\ref{sec:xai} analyses the
Grad-CAM++ and LIME explainability panels. Section~\ref{sec:discussion}
discusses the findings, compares them to prior multi-class systems, and
acknowledges limitations. Section~\ref{sec:conclusion} concludes.
